\documentclass[10pt,a4paper]{amsart}
\usepackage[margin=19mm]{geometry}
\usepackage{amsmath,amssymb,mathtools}
\usepackage[scr=rsfs,cal=euler]{mathalfa}
\usepackage{xfrac}
\usepackage[unicode,hidelinks,bookmarks=false]{hyperref}
\newcommand{\ie}{i.e.\ }
\newcommand{\cf}{cf.\ }
\newcommand{\resp}{respectively\ }
\newcommand{\Subsection}{\subsection}
\providecommand{\nobreakdash}{\nobreak\hbox{-}}
\setlength{\emergencystretch}{2em}
\multlinegap0pt
\usepackage{bm}
\usepackage{bbm}
\usepackage{dsfont}
\usepackage{xspace}
\usepackage{cancel}
\usepackage{mathtools}
\usepackage{amsthm}
\makeatletter
\@ifundefined{thm}{\newtheorem{thm}{Theorem}}{}
\makeatother
\title[Interplay between Contractivity and Monotonicity for Reactio]{Interplay between Contractivity and Monotonicity for Reaction Networks}
\author{Alon Duvall, M. Ali Al-Radhawi, Dhruv D. Jatkar, Eduardo Sontag}
\date{Source: arXiv:2404.18734v2 (CC BY 4.0)}
\begin{document}
\maketitle

\noindent\emph{Source and licence.} Interplay between Contractivity and Monotonicity for Reaction Networks. Version arXiv:2404.18734v2 (CC BY 4.0), distributed under the Creative Commons Attribution 4.0 International licence, as recorded in the arXiv licence metadata for this exact version and shown on the abstract page https://arxiv.org/abs/2404.18734v2. Full rights notice: licenses/arxiv-2404.18734-CC-BY-4.0.txt.\par\smallskip

\noindent\textit{Editorial scope.} Counted unit: the Thm whose source label is \texttt{thm:cubical network cones} in the article (source0.tex lines 1500--1503, source label \texttt{thm:cubical network cones}), delivered together with the article's own complete original proof of it (source0.tex lines 1505--1588, one \texttt{proof} environment of 84 lines). No mathematics is written, repaired or re-derived: statement and proof are the article's own text line for line. Required context retained uncounted: thm:lift viable for monotone (lines 1138--1141). Every definition, display and label of those lines is retained unchanged. Omitted from this package and not redistributed: 1--1137, 1142--1499, 1504--1504, 1589--2317. Nothing inside a retained excerpt is omitted or altered: every retained body is delivered exactly as the article's lines are written, so no omission receipt of any kind accompanies this package. Citations: the retained excerpts carry no citation command at all, so no bibliography entry of the article is transcribed and no reference is redistributed. Reference adaptation: none was needed and none was made; every reference inside the delivered unit resolves inside it, so no unresolved reference or citation is produced. Printed slips kept literal: no printed word, symbol, hypothesis or number of the counted statement or of its proof was altered. Layout and class: the article's own class is \texttt{article} with packages including geometry, url, amsmath, amsthm, amsfonts, amssymb, amscd, mathtools, amsopn, graphicx, epstopdf, svg, float, comment; the wrapper is the lane's pinned \texttt{amsart} editorial wrapper at 10pt on A4, which restates the theorem-like environments the retained text needs and adds the article's own macro definitions that the retained text needs (none), with extra wrapper packages (bm, bbm, dsfont, xspace, cancel, mathtools). The delivered numbering is the unit's own. Assets: no retained span contains a figure environment, a graphics-inclusion command, a PDF-page inclusion, a table or a TikZ picture; no image file is copied into this package, the package declares no asset dependency and no image file is redistributed with it. Licence: version arXiv:2404.18734v2 of this article is distributed under the Creative Commons Attribution 4.0 International licence, as recorded in the arXiv licence metadata for this exact version and shown on the abstract page https://arxiv.org/abs/2404.18734v2. A full rights notice is delivered at licenses/arxiv-2404.18734-CC-BY-4.0.txt. Characters: every retained excerpt is pure ASCII, so the delivered engine, XeLaTeX with its default Latin Modern text font, reports no missing character.
\begin{thm}
\label{thm:lift viable for monotone}
     Suppose we have a viable set of vectors. If we define our cone $K$ to be the lift of this set, then our system is monotone with respect to $K$. 
\end{thm}

\begin{thm}
\label{thm:cubical network cones}
    Cubical networks are monotone with respect to a cubical cone (i.e., a cone with the poset of a cube). 
\end{thm}

\begin{proof}

Define an initial vector \( v \) with coordinates:
    \[
    (v)_i =
    \begin{cases}
        1, & \text{if all reaction vectors have a nonnegative \( i \)th coordinate,} \\
        -1, & \text{if all reaction vectors have a nonpositive \( i \)th coordinate,} \\
        0, & \text{otherwise.}
    \end{cases}
    \]
    Note that if \( v = 0 \), then each row of the stoichiometric matrix must either consist entirely of zeros or contain exactly one \(-1\) and one \(1\). Summing all reaction vectors would then yield the zero vector, contradicting their linear independence. Hence, this procedure ensures that \( v \neq 0 \). By construction, \( v \in Q_2(\Gamma_i) \cup Q_1^+(\Gamma_i) \) for all \( i \).

    Define the set:
    \[
    A = \{ v - \sum_{i \in J}  \Gamma_i \mid J \subset \{1,2,\dots,n\}
     \}.
    \]
    We must show that:
    1. \( A \) forms a cube.
    2. \( A \) is viable. Once we do this, by Theorem \ref{thm:lift viable for monotone} our system is indeed monotone with respect to the lift of this set, which is a cubical cone.

 To confirm the cube structure, consider an invertible transformation \( T \) such that \( T(R_i) = e_i \), where \( e_i \) is the \( i \)th standard basis vector. This maps our set to:
\[
 T(v) - \sum_{i \in J} e_i,
 \]
 which is a cube.


 Note that by construction, \( v \in Q_1(\Gamma_i) \) for all \( i \). Each vector in \( A \) has entries in \( \{0,1,-1\} \). Let \( v_J \in A \) be the vector corresponding to the subset \( J \) in the definition of \( A \). To show that \( A \) is closed, it suffices to verify that  
\[
i \in J \implies v_J \in Q_1(-\Gamma_i), \quad i \notin J \implies v_J \in Q_1(\Gamma_i).
\]

We proceed by case analysis.

\textit{Case 1:} \( i \in J \).  
Each coordinate of \( v_J \) that overlaps with \( \Gamma_i \) must be nonpositive. Consider an arbitrary such coordinate \( l \), and assume without loss of generality that \( (\Gamma_i)_l = -1 \).

If \( \Gamma_i \) is the only reaction involving coordinate \( l \), then \( (v_J)_l = 0 \).

Now suppose there is another reaction \( \Gamma_j \) such that \( (\Gamma_j)_l = -1 \). Since \( (v)_l = -1 \), it follows that:
\[
(v_J)_l =
\begin{cases}
0, & \text{if } j \notin J, \\
1, & \text{if } j \in J.
\end{cases}
\]
Next, suppose instead that \( (\Gamma_j)_l = 1 \). Since \( (v)_l = 0 \), we have:
\[
(v_J)_l =
\begin{cases}
0, & \text{if } j \in J, \\
1, & \text{if } j \notin J.
\end{cases}
\]
In all cases, these conditions ensure that \( v_J \in Q_1(-\Gamma_i) \).

\textit{Case 2:} \( i \notin J \).  
Each coordinate of \( v_J \) that overlaps with \( \Gamma_i \) must be nonnegative. Again, consider an arbitrary such coordinate \( l \) and assume \( (\Gamma_i)_l = -1 \).

If \( \Gamma_i \) is the only reaction involving coordinate \( l \), then \( (v_J)_l = -1 \).

Now suppose there is another reaction \( \Gamma_j \) such that \( (\Gamma_j)_l = -1 \). Since \( (v)_l = -1 \), it follows that:
\[
(v_J)_l =
\begin{cases}
-1, & \text{if } j \notin J, \\
0, & \text{if } j \in J.
\end{cases}
\]
Next, suppose instead that \( (\Gamma_j)_l = 1 \). Since \( (v)_l = 0 \), we have:
\[
(v_J)_l =
\begin{cases}
-1, & \text{if } j \in J, \\
0, & \text{if } j \notin J.
\end{cases}
\]
In all cases, these conditions ensure that \( v_J \in Q_1(\Gamma_i) \).

Since the set is closed under permissible operations, and for each reaction \( \Gamma_i \) there exists some \( k \in A \) such that \( k \in Q_1(\Gamma_i) \cup Q_1(-\Gamma_i) \), we conclude that \( A \) is a viable set.
 \end{proof}

\end{document}
